\documentclass[12pt]{article}
\usepackage[utf8]{inputenc}
\usepackage[T1]{fontenc}
\usepackage{hyperref}
\title{Trinidadian English Creole Language Resources for K-12 Curriculum Development}
\author{Compiled by Claude Sonnet 5 (claude-sonnet-5) for Andr\'as Kornai}
\date{}
\begin{document}
\maketitle
\section{Introduction}
Trinidadian English Creole (ISO 639-3 trf) is an English-based creole of Trinidad and Tobago. Search turned up one major scholarly resource, a dictionary, and little else.
\section{Primary Educational Resources}
I found no primers, textbooks or school programmes in this search.
\section{Linguistic and Reference Materials}
Lise Winer, \emph{Dictionary of the English/Creole of Trinidad \& Tobago: On Historical Principles} (McGill-Queen's University Press, 2009). It is the first scholarly dictionary of the language, with over 12,200 entries giving definitions, spellings, pronunciations, etymologies and usage examples.
\begin{itemize}
\item Publisher page at De Gruyter Brill: \url{https://www.degruyterbrill.com/document/doi/10.1515/9780773576070/html}
\item Listing at 49th Shelf: \url{https://cdlc.49thshelf.com/Books/D/Dictionary-of-the-English-Creole-of-Trinidad-Tobago}
\item Review in Caribbean Beat: \url{https://www.caribbean-beat.com/issue-108/book-reviews-marchapril-2011}
\item Catalogue record (WorldCat): \url{https://search.worldcat.org/oclc/759157066}
\end{itemize}
\section{Religious and Community Resources}
I found none in this search.
\section{Digital and Audio Resources}
I found no free online dictionary or audio resources in this search.
\section{Curriculum Development Context}
The language is not known to be taught in schools. Winer's dictionary is a solid lexical base, but graded teaching materials, a standard orthography for classroom use and audio would all need to be developed.
\end{document}